\documentclass[11pt]{article}
\usepackage[T1]{fontenc}
\usepackage[a4paper,margin=27mm]{geometry}
\usepackage{lmodern,amsmath,amssymb,amsthm,booktabs,microtype,hyperref}
\hypersetup{hidelinks}
\newtheorem{theorem}{Theorem}[section]
\newtheorem{lemma}[theorem]{Lemma}
\newtheorem{proposition}[theorem]{Proposition}
\newcommand{\Q}{\mathbb Q}
\newcommand{\Z}{\mathbb Z}
\newcommand{\F}{\mathbb F}
\newcommand{\N}{\mathrm N}
\newcommand{\Li}{\operatorname{Li}}
\newcommand{\workbenchurl}{https://github.com/KokunoYumeto/zeta-and-l-functions-workbench}
\newcommand{\recordurl}{https://github.com/KokunoYumeto/mathematics-workbenches}
\title{Cubic auxiliary primes and Nullstellensatz consequences,\\
after Adleman--Lenstra and Koiran}
\author{}
\date{24 November 2026}
\begin{document}
\maketitle
\begin{abstract}
Adleman and Lenstra constructed finite fields using small auxiliary primes
under GRH, and Koiran used GRH to place Hilbert's Nullstellensatz in AM.
A uniform zero-free half-plane with boundary strictly below one already
suffices for deterministic construction of cubic finite fields in polynomial
time. We give the conductor-uniform smoothed argument: the named external
seven-eighths theorem yields an auxiliary prime of size
$O((\log(3p))^8)$ for every rational prime $p\ne3$. Full cubic RH would
improve the exponent to two. Liu's main algebraic boundary $B_*$ gives the
smaller exponent $1/(1-B_*)=7.8955174245\ldots$. We distinguish fixed-field and
varying-field Nullstellensatz promises; in particular, the pure cubic
promise is already in AM unconditionally. We also give the precise short
range controlled by smooth squarefree cancellation.
The analytic consequences grant the named external zero-free theorems
and do not prove cubic RH. The mathematics and text were produced by GPT-6 Astra
in ChatGPT Pro mode; the checks identified in the workbench are by
the same model, and no human has checked this text.
\end{abstract}

\section{Main statements and uniform hypotheses}
Put $K=\Q(\omega)$, where $\omega^2+\omega+1=0$.
The named external input is OpenAI (2026a, Theorem~1.1, p.~4):
every finite-order Hecke $L$-function over $K$ has no zero in
$\Re s>7/8$. This theorem from the September 30 release is
granted as an external input. The later October 5 release
(OpenAI, 2026b, Theorem~1.1, p.~2) gives the weaker $11/12$ boundary
and explicitly presents itself as an intermediate step.

More recently, Liu (2026, Theorems~1.2 and 8.29, pp.~3 and 78--80)
proves nonvanishing for every continuous unitary Hecke character of every
number field in the open half-plane with boundary
\[
 B_*=0.873345856106718576801184796065\ldots.                   \tag{1.0}
\]
This is the degree-nine algebraic constant defined by his equations
(42) and (44), pp.~27--29. It is the indicated root in
$(0.8733458561067185768,0.8733458561067185769)$.
There is no conductor or height restriction in that theorem.
We grant this named external input. Its statement permits the
principal pole at one.
His separate one-stage Theorems~1.1 and 8.31 give the simpler rational
value $3493459/4000000$, also displayed below for comparison.
Constants in Liu's analytic estimates may depend on the fixed field
and character. The uniform constants needed here are instead derived
directly from the asserted zero location and classical zero counting.

For a rational prime $p\ne3$, let
\[
 F_p=\Q(\sqrt[3]p),\qquad E_p=K(\sqrt[3]p),\qquad
 L_p(s)=\frac{\zeta_{F_p}(s)}{\zeta(s)}.
\]
The polynomial is Eisenstein at $p$, and its splitting group is $S_3$.
If $\rho_p$ is the standard two-dimensional representation, then
$L_p=L(s,\rho_p)$ is induced from either nontrivial character of
$\operatorname{Gal}(E_p/K)$, and is a nontrivial finite-order Hecke
$L$-function of $K$. It is entire. Its conductor over $\Q$ is
$Q_p=|D_{F_p}|$, and the discriminant of the defining order gives
\[
 Q_p\le27p^2.                                                     \tag{1.1}
\]
Let $\eta=\chi_{-3}$. Artin formalism gives, including the ramified factors,
\[
 \zeta_{E_p}(s)=\zeta(s)L(s,\eta)L_p(s)^2.                          \tag{1.2}
\]
Only degrees one and two occur in the auxiliary-prime argument below.

The Gaussian-period construction in Proposition~\ref{prop:algorithm}
is Adleman and Lenstra's. Theorem~\ref{thm:aux} specifies the uniform
zero-free hypothesis that makes its cubic specialization polynomial.
Koiran's useful-prime gap and arithmetic Nullstellensatz underlie
Theorem~\ref{thm:HN}; its four cases distinguish the required
solution-field hypotheses.


\subsection*{Auxiliary primes and deterministic construction}
\begin{theorem}\label{thm:aux}
Let $p\ne3$ be a rational prime. Suppose $\zeta$, $L(s,\eta)$ and $L_p$ have no nontrivial zeros
with real part greater than a fixed $1/2\le\beta<1$.
There is an effectively computable constant $C_\beta$, independent
of $p$, and an auxiliary prime $\ell\ne3,p$ satisfying
\[
 \ell\equiv1\pmod3,\qquad
 p^{(\ell-1)/3}\not\equiv1\pmod\ell,\qquad
 \ell\le C_\beta\bigl(\log(3p)\bigr)^{1/(1-\beta)}.                 \tag{1.3}\label{eq:aux-prime-bound}
\]
Liu's main theorem therefore gives exponent
$1/(1-B_*)=7.895517424543\ldots$ for every $p\ne3$.
As a corollary using only the named seven-eighths theorem, exponent
eight suffices. RH for the single function $L_p$, together with
the corresponding degree-one RH statements, supplies exponent two
for that input $p$.
\end{theorem}

\begin{proposition}[The deterministic construction]\label{prop:algorithm}
Under OpenAI's September 30 zero-free theorem with boundary $7/8$
for finite-order Hecke characters of $K$, or Liu's theorem with
boundary $B_*$ for unitary Hecke characters over every number field, an
irreducible cubic over $\F_p$ can be constructed deterministically
in time polynomial in $\log p$, for every rational prime $p$.
\end{proposition}

\subsection*{Four quantitative witness promises}
Fix an explicit polynomial $P$. Inputs are sparsely encoded integer
polynomial systems of bit length $N$. A quantitative witness promise
means there is a solution with coordinates in
$\mathcal O_F[1/B]$, where $F$ is a number field,
$B\in\Z\setminus\{0\}$, and
\[
 \log|D_F|,\ \log(2|B|)\le2^{P(N)}.                               \tag{1.4}\label{eq:witness-bound}
\]
No recognition of the promise or computation of $F$ is required.
In part 3 below, a uniform Dirichlet zero-free boundary means one
fixed $1/2\le\beta<1$ such that every nontrivial zero of every
Dirichlet $L$-function, including $\zeta$, has real part at most
$\beta$.

\begin{theorem}\label{thm:HN}
Consider the promise problem with all complex-inconsistent systems
as no-instances and the following systems as yes-instances.
In each case the promise problem is in AM under the stated hypothesis.
\begin{enumerate}
\item If $F$ ranges through a fixed finite collection of number
fields and \eqref{eq:witness-bound} holds, the promise problem is in AM
unconditionally.
\item If $F$ ranges through all pure cubic fields
$\Q(\sqrt[3]a)$ of degree three and \eqref{eq:witness-bound} holds, the promise
problem is in AM unconditionally.
\item If $F$ ranges through all non-Galois cubic fields and
\eqref{eq:witness-bound} holds, a uniform zero-free boundary $\beta<1$ for Dirichlet
$L$-functions suffices.
\item If $F$ ranges through arbitrary number fields and \eqref{eq:witness-bound}
holds, Liu's named all-number-field zero-free theorem suffices.
\end{enumerate}
None of these conclusions asserts AM for arbitrary solution fields from
RH for the single radicand 269.
\end{theorem}

\subsection*{The controlled smooth inverse range}
For this statement the radicand is fixed at 269. Let
$\psi=\psi_{269}$ be either cubic Hecke character of $K$ with
$L_K(s,\psi)=L_{269}(s)$, and let $S$ be the primes above 3 and 269.
All ideals below are coprime to $S$. Put
\[
 \begin{gathered}
 P(s)=L_K^S(s,\psi),\\
 b(\mathfrak c)=\mu_K(\mathfrak c)^2\psi(\mathfrak c),\qquad
 \ell(\mathfrak d)=(-1)^{\Omega(\mathfrak d)}\psi(\mathfrak d),
 \end{gathered}
\]
where $\Omega$ counts prime factors with multiplicity. The last
sequence is the convolution inverse of $b$, and the Dirichlet
series of $b$ is $P(s)/P(2s)$. Assume that $P$ has no zero in
$\Re s>\beta$, where $1/2\le\beta<1$ is fixed. Then, for every
$V\in C_c^\infty((0,\infty))$ and every $\epsilon>0$,
\[
 \sum_{\mathfrak c}b(\mathfrak c)V(N\mathfrak c/Y)
       =O_{V,\beta,\epsilon}(Y^{\beta/2+\epsilon})\qquad(Y\ge1).
\]
The character is fixed, so these constants need not be uniform
in varying radicands or conductors.
\begin{proposition}\label{prop:short}
With $P$, $b$, $\ell$, $S$ and $\beta$ as just defined,
for fixed $W\in C_c^\infty((1,2))$, define
\[
 T_{\le D}(X)=
 \sum_{\N\mathfrak d\le D}\ell(\mathfrak d)
 \sum_{\mathfrak c}b(\mathfrak c)\log\N\mathfrak c\,
          W(\N\mathfrak c\N\mathfrak d/X).
\]
For $1\le D\le X$ and every $\epsilon>0$,
\[
 T_{\le D}(X)\ll_{W,\beta,\epsilon}
       X^{\beta/2+\epsilon}D^{1-\beta/2}\log(2X).
                                                                    \tag{1.5}\label{eq:short-inverse-bound}
\]
In particular $D=X^{(1-\beta)/(2-\beta)}$ is admissible at exponent
$1/2+\epsilon$.
\end{proposition}

At $\beta=B_*$ the cutoff exponent is
$(1-B_*)/(2-B_*)=0.112416170108\ldots$.
The same Mellin shift gives exponent $B_*/2+\epsilon$ for each
fixed nonnegative integer power $P(s)^j/P(2s)$, including
$8P(s)^4/P(2s)$ with coefficients
$8\psi(\mathfrak a)\tau_K(\mathfrak a)^2$, where $\tau_K$ counts
ideal divisors.
The complementary inverse range
$X^{(1-B_*)/(2-B_*)}<N\mathfrak d<2X$ is not bounded here at
the square-root scale, and no quarter-power coefficient bound
is asserted.

\subsection*{A conditional fixed-field prime count}
For $E=\Q(\omega,\sqrt[3]{269})$, let $\pi_C(x)$ count
rational primes at most $x$ that are unramified in $E$ and have
Frobenius in a conjugacy class $C\subset S_3$; set
$\delta_C=|C|/6$ and $\Li(x)=\int_2^xdu/\log u$.
Under RH for $L_{269}$, $\zeta$ and $L(s,\eta)$,
Greni\'e and Molteni (2019, Corollary~1.2) give, for every $x\ge2$,
\begin{align*}
 |\pi_C(x)-\delta_C\Li(x)|&\le\delta_C\sqrt x\,\mathcal E(x),\\
 \mathcal E(x)&=
 \left(\frac1{2\pi}+\frac3{\log x}\right)\log(3^3\,269^4)\\
 &\quad+6\left(\frac{\log x}{8\pi}+\frac1{4\pi}+\frac6{\log x}\right).
\end{align*}
The three-cycle, transposition and identity classes have densities
$1/3$, $1/2$ and $1/6$, respectively. Requiring additional congruences
modulo a varying integer introduces further twist RH hypotheses;
they do not follow from RH for $L_{269}$ alone.

\section*{Proof overview}
The main analytic issue is uniformity in the conductor. A sharp
prime-counting error often introduces an extra logarithm of the range;
with a fixed smooth weight the zero sum is instead bounded by the
logarithm of the conductor alone. The Hadamard product gives the needed
local zero count with an absolute constant. A zero-free boundary
$\beta<1$ then makes the three-cycle prime sum positive as soon as
$x^{1-\beta}$ exceeds a constant times $\log p$.
Subtracting higher prime powers and the two ramified primes is
uniform in $p$.

The auxiliary prime makes Frobenius a three-cycle on three Gaussian
periods; their invariant algebra is consequently the desired cubic
field. For Nullstellensatz, degree-one prime ideals reduce a bounded
algebraic solution to many finite-field solutions. Koiran's unconditional
bound on useful primes for an inconsistent system supplies the opposite
case, and a pairwise-independent hash gives the AM protocol.
The final section treats prime extraction separately: the available
smooth squarefree estimate controls a short convolution range, but
does not bound its complementary range at the square-root scale.

\subsection*{The conductor-uniform explicit formula}
We establish the uniform estimate used in both applications.

\begin{lemma}[Uniform smooth explicit formula]\label{lem:uniform}
Fix $W\in C_c^\infty((1,2))$. For one of $\zeta(s)$, $L(s,\eta)$,
or $L_p(s)$, denote its logarithmic-derivative coefficients by
$\Lambda_H(n)$, its conductor by $Q_H$, and its pole order at one
by $\delta_H$. Suppose every nontrivial zero has real part at most
$\beta$, where $1/2\le\beta<1$ is fixed. Then, uniformly in $p$,
\[
 \sum_n\Lambda_H(n)W(n/x)
 =\delta_H x\widehat W(1)
       +O_{W,\beta}\bigl(x^\beta\log(3Q_H)\bigr),\qquad x\ge2.      \tag{1.6}\label{eq:uniform-explicit-formula}
\]
Here $\widehat W(s)=\int_0^\infty W(u)u^{s-1}\,du$.
The constant is effective and depends on the fixed weight and $\beta$,
not on the character or its conductor.
\end{lemma}

\begin{proof}
These are primitive $L$-functions of degree at most two with one of a
fixed finite set of archimedean factors. Their standard explicit
formula, obtained by Mellin inversion of $-H'/H$, is
\[
 \delta_H x\widehat W(1)
 -\sum_{\rho}m_\rho x^\rho\widehat W(\rho)
 -\sum_{\rho_{\rm triv}}m_{\rho_{\rm triv}}
             x^{\rho_{\rm triv}}\widehat W(\rho_{\rm triv}).
                                                                    \tag{1.7}\label{eq:uniform-zero-expansion}
\]
The possible zero at zero belongs to the last sum. The functional
equations and the same upper bound for the conjugate character put
nontrivial zeros in $1-\beta\le\Re\rho\le\beta$.
For $L_p$ the coefficients are real, so the conjugate function is the
same one; the degree-one factors here are real as well.
Repeated integration by parts gives, for every $A$,
\[
 |\widehat W(\sigma+it)|\le C_{W,A}(1+|t|)^{-A}
       \quad(0\le\sigma\le1).
 \]
The conductor-uniform local zero count is
\[
 \#\{\rho:j\le|\Im\rho|<j+1\}
 \ll \log\bigl(3Q_H(j+3)^2\bigr),\qquad j\ge0.                     \tag{1.8}\label{eq:uniform-zero-count}
\]
Here is the uniformity argument. Complete $H$ with its conductor
and gamma factors, and multiply by $s(s-1)$ in the zeta case;
call the resulting entire function $\mathcal H$. Its order-one
Hadamard product, reality, and functional equation give
\[
 \Re\frac{\mathcal H'}{\mathcal H}(2+it)
 =\sum_\rho m_\rho
   \frac{2-\Re\rho}{(2-\Re\rho)^2+(t-\Im\rho)^2}.
\]
For completeness, the real part of the canonical-product
logarithmic derivative is a real constant plus
$\sum_\rho\Re(1/(s-\rho))$. The latter series converges absolutely,
since the zeros lie in $0\le\Re\rho\le1$ and the function has order
one. Both the logarithmic derivative and this series change sign
under $s\mapsto1-\overline s$, by the functional equation and the
zero symmetry. The constant is therefore zero.
Each zero with $|t-\Im\rho|\le1$ contributes
at least $1/5$, since $0\le\Re\rho\le1$.
On the other hand, the Euler product bounds $H'/H(2+it)$ by an
absolute constant: its local roots have modulus at most one
and its degree is at most two. The logarithmic derivatives
of the fixed gamma factors are $O(\log(|t|+3))$, the conductor
term is $\tfrac12\log Q_H$, and the zeta pole correction is
$O(1)$. This proves \eqref{eq:uniform-zero-count}. In particular, no constant depending
arbitrarily on $p$ is introduced. Taking $A=3$
in \eqref{eq:uniform-zero-count} bounds the nontrivial-zero sum in absolute value by
$O_W(x^\beta\log(3Q_H))$, including zeros on $\Re s=\beta$.
The trivial zeros are nonpositive integers, with multiplicities at
most two. Since $W$ is supported in a fixed compact subinterval of
$(1,2)$, their series is $O_W(1)$ for $x\ge2$.
The Mellin contour proof of \eqref{eq:uniform-zero-expansion} is justified on good horizontal
heights by the same local zero count; its zero series converges
absolutely. This proves \eqref{eq:uniform-explicit-formula}, with no loss $x^\epsilon$.
The completed function, order-one product, and logarithmic derivative
are written explicitly in Greni\'e and Molteni (2019, (2.5)--(2.8)
and (3.1)); their unconditional Appendix~A, Theorem~A.1, also gives
a conductor-uniform zero count. The Mellin-contour formula is the
classical explicit formula; see also Iwaniec and Kowalski
(2004, Chapter~5).
\end{proof}

\section{Auxiliary primes and cubic finite fields}
We apply the uniform estimate to the three-cycle class and then give
the finite-field construction it permits.
For an unramified prime $\ell\nmid3p$, let $c_p(\ell)$ be the
indicator that its Frobenius in $S_3$ is a three-cycle. The
character table gives the exact identity
\[
 c_p(\ell)=\frac{1+\eta(\ell)-\operatorname{tr}\rho_p(\ell)}3
 =\mathbf1_{\{\ell\equiv1\ (3),\ p^{(\ell-1)/3}\not\equiv1\ (\ell)\}}.
                                                                    \tag{2.1}
\]
The power test is applied after excluding $\ell=p$. For example, the first
formula has values $0,0,1$ at the identity, a transposition, and
a three-cycle, respectively.

\begin{proof}[Proof of Theorem~\ref{thm:aux}]
Fix a nonzero nonnegative $W\in C_c^\infty((1,2))$ once and for all.
Apply Lemma~\ref{lem:uniform} to the character combination
$(1+\eta-\rho_p)/3$. Its prime-power coefficients away from $3p$
are $c_p(\operatorname{Frob}_\ell^k)\log\ell$, each between zero
and $\log\ell$. Thus the prime-power terms of exponent at least two
have absolute total at most
\[
 \|W\|_\infty\sum_{k=2}^{\lfloor\log_2(2x)\rfloor}
                  \vartheta((2x)^{1/k})\ll_W\sqrt x.              \tag{2.3}
\]
Chebyshev's elementary bound handles $k=2$. The other terms have
total $O(x^{1/3}\log(2x))=O(\sqrt x)$, uniformly
for $x\ge2$. At either ramified prime, each local trace is bounded
by its representation degree. The entire contribution of its
powers in the support of $W(n/x)$ is $O_W(\log(2x))$: if a power
occurs its prime is at most $2x$, and the number of exponents
times its logarithm is at most $\log(2x)$. Deleting both ramified
primes costs $O_W(\log(2x))$, with no dependence on $p$.
Consequently
\[
 \sum_{\substack{\ell\ {\rm prime}\\\ell\nmid3p}}
          c_p(\ell)\log\ell\,W(\ell/x)
 =\frac{x\widehat W(1)}3+
        O_{W,\beta}\bigl(x^\beta\log(3p)\bigr).                   \tag{2.4}
\]
Take $x=A(\log(3p))^{1/(1-\beta)}$.
The ratio of the error to the main term is at most a fixed
constant times $A^{\beta-1}$. An effectively computable fixed
$A$ makes this less than one. The sum is then positive, and a
prime in its support lies between $x$ and $2x$ and satisfies
(2.1). This proves \eqref{eq:aux-prime-bound}.

Either named external theorem applies to $L_p$ and the
trivial Hecke function $\zeta_K=\zeta L(\eta)$, whose factorization
supplies the required zero-free half-plane for the degree-one
factors away from their pole at one. Nonvanishing there is
classical. Alternatively the stipulated Dirichlet RH supplies
these two factors. No extra conjectural twist family is needed.
\end{proof}

\begin{proof}[Proof of Proposition~\ref{prop:algorithm}]
For $p=3$, use $X^3-X-1$. For $p\ne3$, enumerate rational primes
$\ell$ and test the two congruences in \eqref{eq:aux-prime-bound}, explicitly rejecting
$\ell=p$. Trial division suffices for primality and factorization
of $\ell-1$: the theorem bounds every searched integer by a
polynomial in $\log p$. The search and modular exponentiations
therefore take polynomial time.

Let $H$ be the subgroup of cubes in $(\Z/\ell\Z)^\times$ and form
the finite etale algebra
\[
 A=\F_p[T]/\Phi_\ell(T),\qquad B=A^H,
 \quad h:T\longmapsto T^h .
\]
The polynomial $\Phi_\ell$ is squarefree because $p\ne\ell$.
Taking invariants commutes with a field extension because it is
the common kernel of linear maps $h-1$. After passage to an
algebraic closure, $B$ has three points,
indexed by the three cosets of $H$. Frobenius acts by multiplication
by $p$ on those cosets, a three-cycle by \eqref{eq:aux-prime-bound}. Hence $B$ is the
field $\F_{p^3}$.
The period $\xi=\sum_{h\in H}T^h$ belongs to $B$ and is not scalar:
otherwise its polynomial, minus a constant, would be a scalar
multiple of $\Phi_\ell$, whereas its nonconstant coefficients
take both values zero and one. Since $[B:\F_p]=3$, $\xi$ has
degree three. Compute multiplication by $\xi$ on $B$ by linear
algebra in the algebra of dimension $\ell-1$; its characteristic
polynomial is the required irreducible cubic. All operations
are polynomial in $\ell$ and $\log p$.
This is the Gaussian-period construction of Adleman and
Lenstra (1986, Algorithm~A and Lemma~6), with its prime-selection
input supplied by (2.4).
\end{proof}

For the single input $p=269$, no analytic input is needed:
$\ell=7$ works because $269^2\equiv2\pmod7$. The corresponding
period polynomial is $X^3+X^2-2X-1$, irreducible modulo 269.
The accompanying exact program verifies this assertion.
Thus RH for this one fixed $L$-function improves quantitative
prime distribution; it is not what makes this one finite field
constructible. A proof only at $p=269$ supplies no uniform
algorithmic premise for all other radicands.

\subsection{Exactly what fixed-269 RH gives for prime distribution}
Let $E=\Q(\omega,\sqrt[3]{269})$ and let $\pi_C(x)$ count rational
primes at most $x$, unramified in $E$, with Frobenius in a conjugacy
class $C\subset S_3$. Put $\delta_C=|C|/6$ and
$\Li(x)=\int_2^xdu/\log u$. Assume RH for $L_{269}$, $\zeta$, and
$L(s,\eta)$. Equation (1.2) then implies RH for $\zeta_E$.
The cubic discriminant is $-3\cdot269^2$: at 269 the different is
tame of exponent two; at three, 269 is a cube in $\Q_3$ since
$-269\in1+9\Z_3$ and cubing maps $1+3\Z_3$ onto $1+9\Z_3$ by
the three-adic logarithm. The local algebra is
$\Q_3\times\Q_3(\omega)$. Thus its three-adic discriminant exponent
is one. All other primes are excluded by the polynomial discriminant
$-27\cdot269^2$. The conductor-discriminant formula applied to (1.2) gives
\[
 |D_E|=3(3\cdot269^2)^2=3^3\,269^4.
\]
Greni\'e and Molteni's Corollary~1.2 therefore gives, for every $x\ge2$,
\begin{align*}
 |\pi_C(x)-\delta_C\Li(x)|
 &\le \delta_C\sqrt x\,\mathcal E(x),\\
 \mathcal E(x)
 &=\left(\frac1{2\pi}+\frac3{\log x}\right)\log(3^3\,269^4)\\
 &\quad+6\left(\frac{\log x}{8\pi}
              +\frac1{4\pi}+\frac6{\log x}\right).
                                                               \tag{2.5}
\end{align*}
Their definition of $\pi_C$ excludes ramified primes, as ours does.
The three interpretations are
\[
\begin{array}{c|c|c}
 C&\delta_C&\text{condition at }\ell\nmid3\cdot269\\ \hline
 1&1/6&\ell\equiv1\ (3),\quad269\text{ a cube modulo }\ell\\
 \text{transpositions}&1/2&\ell\equiv2\ (3)\\
 \text{three-cycles}&1/3&\ell\equiv1\ (3),\quad269\text{ not a cube modulo }\ell .
\end{array}
\]
This is a direct conditional specialization of their explicit theorem.
Requiring an additional residue class modulo a varying integer $m$
introduces the twists $L(s,\rho_{269}\otimes\chi)$ for characters
$\chi\bmod m$, together with degree-one factors.
Apart from twists isomorphic to the original representation, RH for
$L_{269}$ does not furnish these additional hypotheses.

\section{Precisely scoped Nullstellensatz consequences}
We specify the solution-field promise before drawing the complexity
conclusion, so that no varying-field assertion is inferred from a
single fixed character.
The quantitative witness promise and its four cases are stated
in Section~1, in Theorem~\ref{thm:HN}.

\begin{proof}[Proof of Theorem~\ref{thm:HN}]
For part 1 the unconditional prime ideal theorem in each fixed field
gives $\pi_F(x)\sim\Li(x)$. Prime ideals of residue degree at least
two contribute at most $[F:\Q]\pi(\sqrt x)$. Deleting prime ideals
above divisors of $B$ costs at most $[F:\Q]\omega(B)$.
Each remaining degree-one prime ideal reduces the promised solution
to $\F_\ell$; a rational prime is counted at most $[F:\Q]$ times.
For this fixed finite collection, the number of useful rational
primes up to $x$ is therefore at least
$c x/\log x-\omega(B)$ once $x$ exceeds a fixed threshold.

For part 2, the quadratic resolvent is always $\Q(\omega)$.
Every unramified prime $\ell\equiv2\pmod3$ has transposition
Frobenius in $S_3$ and therefore a degree-one prime ideal in $F$.
The fixed-modulus prime number theorem gives
$\pi(x;3,2)\sim\Li(x)/2$, with constants independent of $a$.
Deleting the divisors of $D_FB$ costs at most
$\log_2(2|D_FB|)\le O(2^{P(N)})$ rational primes.
Thus a sufficiently large $x=2^{Q(N)}$, for a fixed polynomial
$Q$, supplies the required useful primes unconditionally.
Equivalently, the cube map on $\F_\ell^\times$ is bijective at
these primes; the Frobenius argument avoids any choice of a
possibly nonmaximal cubic order.

For part 3, let $\Q(\sqrt d)$ be the quadratic resolvent, with
fundamental discriminant $d$. Then $|d|\le|D_F|$.
The condition $\chi_d(\ell)=-1$ again forces transposition
Frobenius, hence a degree-one prime ideal.
The proof of Lemma~\ref{lem:uniform} applies uniformly to every
primitive quadratic Dirichlet character: only the two fixed
degree-one gamma factors occur. Applying it to
$(1-\chi_d)/2$, and deleting prime powers and divisors of $D_FB$,
gives a positive multiple of $x/\log x$ useful primes whenever
$x^{1-\beta}\ge C\log(3|D_F|)$ and the denominator loss is smaller
than that main term. Both requirements hold for a sufficiently
large fixed polynomial $Q$ in $x=2^{Q(N)}$.
This transposition argument is also used in the note contributed by
Bhattacharya (2026, proof of Theorem~17.1, pp.~56--57), whose
title page credits GPT-6 Astra, for cubic class-group torsion; its use here
is in Koiran's useful-prime criterion.

For part 4 put $n_F=[F:\Q]$. The Hadamard proof in
Lemma~\ref{lem:uniform}, now with $n_F$ gamma factors, gives
\[
 \sum_{\mathfrak p,k}\log\N\mathfrak p\,
       W((\N\mathfrak p)^k/x)
 =x\widehat W(1)+
 O_W\bigl(x^{B_*}(\log(3|D_F|)+n_F)\bigr).                   \tag{3.2}
\]
The constant is independent of the field: the local zero count is
$O(\log(3|D_F|)+n_F\log(|t|+3))$, the Euler logarithmic derivative
at two is $O(n_F)$, and the trivial-zero multiplicities are at
most $n_F$. Higher prime powers and prime ideals of residue
degree at least two cost $O_W(n_F\sqrt x)$.
Minkowski's discriminant bound gives
\[
 |D_F|\ge
 \left(\frac{\pi}{4}\right)^{2r_2}
 \left(\frac{n_F^{n_F}}{n_F!}\right)^2
 \ge\frac12\left(\frac{\pi}{2}\right)^{n_F},
\]
where $r_2$ is the number of complex pairs and the last step uses
$n_F!\le n_F^{n_F}2^{-\lfloor n_F/2\rfloor}$.
Thus $n_F\ll\log(2|D_F|)$; the classical input is recorded,
for example, in Milne (2020, Theorem~4.9 and its proof, p.~72).
Once $x^{1-B_*}\ge C_W\log(3|D_F|)$, division by the maximum
$n_F$ prime ideals above a rational prime gives at least
$c_Wx/(n_F\log x)-\omega(B)$ useful rational primes.
In view of \eqref{eq:witness-bound}, taking $x=2^{Q(N)}$ for a sufficiently large
fixed polynomial $Q$ meets all these inequalities.
Every selected prime still has polynomial bit length.

Here is the public-coin protocol, including its soundness.
Koiran's unconditional Theorem~5 bounds the number of useful
primes for an inconsistent system by $M=2^{q(N)}$, for a fixed
explicit polynomial $q$. It follows from an arithmetic
Nullstellensatz identity $\sum g_i f_i=b\ne0$:
every useful prime divides $b$, with $\log_2|b|\le2^{q(N)}$.
Increase $Q$ so that every promised yes-instance has at least
$12M$ useful primes in the selected polynomial-bit universe.
Choose $m$ with $3M\le2^m<6M$, and choose a uniformly random
affine binary map $h:\F_2^k\to\F_2^m$, where $k$ encodes that universe.
Merlin supplies a prime $\ell$ encoded in $k$ bits with $h(\ell)=0$, a primality
certificate, and a solution modulo $\ell$. The verifier checks
all of them in polynomial time; sparse powers are evaluated by
repeated squaring.

For a no-instance, the union bound gives acceptance probability
at most $M/2^m\le1/3$. For a yes-instance, the pairwise-independent
indicators of $h(\ell)=0$ have total expectation
$\mu\ge12M/2^m>2$ and variance at most $\mu$.
The second-moment inequality gives probability at least
$\mu/(1+\mu)>2/3$ of a valid witness.
Primality has polynomial certificates, or can be tested
deterministically. This proves AM.
Koiran's DIMACS Theorem~2 supplies the original complexity
framework; the alternative single-hash proof here uses the
same useful-prime gap. His 1997 erratum corrects Theorem~8;
it does not affect the unconditional Theorem~5 used above.
\end{proof}

In particular, a bounded-witness promise with solution fields among
the subfields of the fixed splitting field
$\Q(\omega,\sqrt[3]{269})$ is already covered by part 1.
A varying family of all non-Galois cubic solution fields includes
different quadratic resolvents; part 3 treats it through Dirichlet
characters rather than cubic Hecke factors.
Part 4 explicitly grants a theorem over every number field.
No such family hypothesis is inferred from the fixed character 269.

\section{The controlled smooth inverse range}
We quantify the part of prime extraction controlled by the available
squarefree estimate.
Let $\psi=\psi_{269}$ be either cubic Hecke character with
$L_K(s,\psi)=L_{269}(s)$, and let $S$ consist of the prime ideals
above 3 and 269. All ideals in the sums below are coprime to $S$.
Since conjugation inverts $\psi$ and preserves $S$,
$L_K^S(s,\psi^2)=L_K^S(s,\psi)$.
The Euler products therefore give
\[
 P(s)=L_K^{S}(s,\psi),\qquad
 B(s)=\frac{P(s)}{P(2s)}
       =\sum_{\mathfrak c}\frac{b(\mathfrak c)}{(\N\mathfrak c)^s},
 \quad b(\mathfrak c)=\mu_K(\mathfrak c)^2\psi(\mathfrak c).
\]
Its convolution inverse is
$\ell(\mathfrak d)=(-1)^{\Omega(\mathfrak d)}\psi(\mathfrak d)$,
where $\Omega$ counts prime factors with multiplicity.
Suppose $P$ has no zero in $\Re s>\beta$, where $1/2\le\beta<1$.
Then, for each
$V\in C_c^\infty((0,\infty))$ and $\epsilon>0$,
\[
 \sum_{\mathfrak c}b(\mathfrak c)V(\N\mathfrak c/Y)
       =O_{V,\beta,\epsilon}(Y^{\beta/2+\epsilon}),\qquad Y\ge1.    \tag{4.1}
\]
by the zero-free denominator $P(2s)$ and a Mellin shift.
Here are the growth details. On any closed inner strip of
$\Re z>\beta$, take an analytic logarithm of $P(z)$, with its
branch fixed by the absolutely convergent Euler logarithm at
$2+it$. A disk centered there and lying in a slightly larger
inner strip reaches any prescribed point of the smaller strip.
The standard functional equation and convexity bound give
$\log|P(z)|=O(\log(2+|t|))$ from above on that disk, whereas
$\log P(2+it)=O(1)$. Borel--Carath\'eodory therefore gives
$|\log P(z)|=O(\log(2+|t|))$ on the smaller disk.
Consequently $1/P(z)$ has polynomial growth on each such strip.
The finitely many deleted Euler factors are uniformly bounded
away from zero there and change none of these conclusions.

The numerator $P(s)$ also has polynomial growth in closed vertical
strips. Mellin inversion for the left-hand side of (4.1) is thus
shifted from $\Re s>1$ to $\Re s=\beta/2+\epsilon$ without poles:
$P(2s)$ has no zero there, and the Mellin transform of $V$
decays faster than any power. This proves (4.1), with constants
depending on the fixed character, weight, and $\epsilon$.

\begin{proof}[Proof of Proposition~\ref{prop:short}]
For $Y=X/\N\mathfrak d$, the inner sum is
$\log Y$ times (4.1) with weight $W$, plus (4.1) with
weight $u\mapsto W(u)\log u$. Both weights are fixed, independently
of $\mathfrak d$. Hence its absolute value is
$O_{W,\beta,\epsilon}(Y^{\beta/2+\epsilon}\log(2X))$.
Now $|\ell|\le1$ and the number of ideals of norm at most $t$
is $O(t)$. Partial summation gives
$\sum_{\N\mathfrak d\le D}(\N\mathfrak d)^{-\beta/2}
=O_\beta(D^{1-\beta/2})$.
Finally $\beta/2+(1-\beta/2)(1-\beta)/(2-\beta)=1/2$.
\end{proof}

Thus the named external bounds give the exact comparison
\[
\begin{array}{c|c|c}
 \text{zero-free boundary}&
 \text{squarefree exponent}&\text{exponent in the cutoff }D\\ \hline
 7/8&7/16&1/9\\
 3493459/4000000&3493459/8000000&506541/4506541\\
 B_*&B_*/2&(1-B_*)/(2-B_*) .
\end{array}
\]
The strongest cutoff exponent is $0.112416170108\ldots$, larger than
the rational intermediate value $0.112401285154\ldots$ and $1/9$. The same Mellin shift applies to every fixed nonnegative integer power
$P(s)^j/P(2s)$. In particular the Euler product
$8P(s)^4/P(2s)$ has coefficients
$8\psi(\mathfrak a)\tau_K(\mathfrak a)^2$, where $\tau_K$ counts
ideal divisors: this follows locally from
$\sum_{k\ge0}(k+1)^2z^k=(1-z^2)/(1-z)^4$.
Its smooth coefficient sum therefore has exponent
$B_*/2+\epsilon=0.436672928053359\ldots+\epsilon$ as well. No quarter-power bound
is proved here.

\subsection*{What remains open}
The complete convolution is the coefficient sum of $-B'/B$;
its prime terms are $\psi(\mathfrak p)\log\N\mathfrak p$.
Prime powers cost $O_W(\sqrt X)$.
The improvement in \eqref{eq:short-inverse-bound} controls a larger range than the
polylogarithmic cutoff from the elementary square-root
squarefree sum. It does not estimate the remaining range.
Liu's main theorem bounds the complete prime sum
at exponent $B_*$, so it also bounds that remaining range at
$B_*+\epsilon$ after subtracting \eqref{eq:short-inverse-bound}. No square-root
bound for the long range
$X^{(1-B_*)/(2-B_*)}<\N\mathfrak d<2X$ has been established.

\section*{Note on method}
The mathematics and text were produced by GPT-6 Astra in ChatGPT
Pro mode. The original program \emph{Cubic follow-up certificate}
uses exact integer arithmetic for the character-table identity and
the Gaussian-period examples. The workbench identifies the scope
of the author's checks and the separate checks by the same model;
no human has checked this text. The text and original program are
dedicated to the public domain under CC0 1.0. Programs, statements,
and check records are at \url{\workbenchurl}; permanent version
records are at \url{\recordurl}.

\begin{thebibliography}{99}
\small
\setlength{\itemsep}{3pt}
\bibitem{OpenAI1} OpenAI. (2026a, September 30).
\emph{The quasi-Riemann hypothesis: A zero-free half-plane
$\Re s>7/8$}. Theorem 1.1, p.~4.
\href{https://github.com/openai/math/blob/main/preprints/The-Quasi-Riemann-Hypothesis-September-30-2026/paper.pdf}{Dated primary release}.
\bibitem{OpenAI2} OpenAI. (2026b, October 5).
\emph{The quasi-Riemann hypothesis}. Theorem 1.1, p.~2.
\href{https://github.com/openai/math/blob/main/preprints/The-Quasi-Riemann-Hypothesis-October-5-2026/paper2.pdf}{Dated primary release}.
\bibitem{Liu} Liu, B. (2026, November 13).
\emph{Joint marked moments and optimized zero-free half-planes
for unitary Hecke $L$-functions} (Version 2).
Theorems 1.1--1.2, pp.~2--3; equations (42), (44), pp.~27--29;
Theorem 8.29, pp.~78--80; Theorem 8.31, pp.~80--83.
\url{https://hexagonmath.org/2611.00421v2}.
\bibitem{IK} Iwaniec, H., \& Kowalski, E. (2004).
\emph{Analytic number theory}. American Mathematical Society.
Chapter 5, completed $L$-functions, zero counts and explicit formulae.
\bibitem{GM} Greni\'e, L., \& Molteni, G. (2019).
\emph{An effective Chebotarev density theorem under GRH}
(Version 3, May 28). Corollary 1.2; equations (2.5)--(2.8) and
(3.1); Appendix A, Theorem A.1.
\url{https://arxiv.org/abs/1709.07609v3}.
\bibitem{Bhattacharya} Bhattacharya, S. (2026, October 11).
\emph{Unconditional implications of the Quasi-Riemann Hypothesis}.
Contributed by S. Bhattacharya; title-page attribution: GPT-6 Astra,
manuscript dated October 9. Proof of Theorem 17.1, pp.~56--57,
the non-Galois cubic transposition argument.
\url{https://hexagonmath.org/2610.00178v1}.
\bibitem{MilneANT} Milne, J. S. (2020).
\emph{Algebraic number theory} (Version 3.08).
Theorems 4.3 and 4.9, pp.~70--72.
\url{https://www.jmilne.org/math/CourseNotes/ANT.pdf}.
\bibitem{AL} Adleman, L. M., \& Lenstra, H. W., Jr. (1986).
Finding irreducible polynomials over finite fields.
\emph{Proceedings of STOC 1986}, 350--355.
Algorithm A, pp.~350--351; Lemma 6, p.~352.
\url{https://pub.math.leidenuniv.nl/~lenstrahw/PUBLICATIONS/1986a/art.pdf}.
\bibitem{Koiran} Koiran, P. (1996).
\emph{Hilbert's Nullstellensatz is in the polynomial hierarchy}.
DIMACS Technical Report 96-27.
Theorem 2, printed pp.~5--6; Theorem 5, printed p.~9.
\href{https://perso.ens-lyon.fr/pascal.koiran/Publis/DIMACS.96-27.ps}{Author's extended report}.
\bibitem{Erratum} Koiran, P. (1997, August 8).
\emph{Erratum to Hilbert's Nullstellensatz is in the polynomial hierarchy}.
\href{https://perso.ens-lyon.fr/pascal.koiran/Publis/hn.erratum.ps}{Author's erratum}, p.~1.
\end{thebibliography}
\end{document}
